%! Solving Exponential Equations with a Common Base — Unit 7, Lesson 7.3 — Student Packet Cover Sheet
\documentclass[10pt]{article}
\usepackage{algebra2-article}
\usepackage{algebra2-boxes}

\begin{document}

% Full-bleed forest banner
\begin{tikzpicture}[remember picture, overlay]
  \fill[forest] (current page.north west)
    rectangle ([yshift=-0.9in]current page.north east);
\end{tikzpicture}
\vspace{-0.8in}

\begin{center}
  {\color{white}\textbf{\LARGE Algebra 2: Shepherd}}\\[4pt]
  {\color{forestmid}\large\textbf{Unit 7 \quad Exponential Functions}}\\[6pt]
  {\color{white}\textbf{\large Lesson 7.3 \quad Solving Exponential Equations with a Common Base}}
\end{center}

\vspace{0.22in}
\namedateperiod
\vspace{0.18in}

\begin{learningtargetbox}
\small
\begin{enumerate}[leftmargin=1.5em, itemsep=5pt]
  \item \ldots{} solve an exponential equation by rewriting \textbf{both sides as powers of one base} and then setting the exponents equal --- including cases that need Unit 6's negative and fractional exponents.
  \item \ldots{} explain \textbf{why} that move is allowed: an exponential function never produces the same output twice, so equal bases force equal exponents.
  \item \ldots{} find a solution on a graph as the crossing point of a curve and a horizontal line, and use the \textbf{range} to decide when an exponential equation has no solution at all.
\end{enumerate}
\end{learningtargetbox}

\vspace{0.14in}

\begin{tocbox}
\small
\renewcommand{\arraystretch}{1.5}
\begin{tabularx}{\linewidth}{@{} c l X r @{}}
\rowcolor{forestbg}
\textbf{\#} & \textbf{Component} & \textbf{Description} & \textbf{Score} \\
\midrule
1 & Warm-Up        & Spiral review: rewriting numbers as powers, add vs.\ multiply, and a question about shape & \blank{1.2cm} \\
2 & Guided Notes   & The one-to-one property, matching bases, the three steps, and solutions as intersections & \blank{1.2cm} \\
3 & Group Activity & \emph{Make the Bases Match} --- rewriting, harder equations, error analysis, and a colony, by tier & \blank{1.2cm} \\
4 & Exit Ticket    & Quick check: two equations solved, one explained from the range, and one that cannot be solved yet & \blank{1.2cm} \\
5 & Homework       & Eight equations, a graph read both directions, a sorting task, error analysis, and an equation with no solution & \blank{1.2cm} \\
\midrule
  & \multicolumn{2}{r}{\textbf{Total}} & \blank{1.2cm} \\
\end{tabularx}
\end{tocbox}

\vspace{0.10in}

\begin{remindbox}
\small An \textbf{exponential equation} puts the variable in the exponent, and today you solve one for
the first time. The whole method is three steps: rewrite \textbf{both} sides as powers of a single
base, set the exponents equal, and solve what is left. Step 1 is Unit 6 work --- $8=2^{3}$,
$\frac{1}{27}=3^{-3}$, $\sqrt{2}=2^{1/2}$ --- and it is where nearly every mistake happens. Watch the
two exponent rules that look alike: raising a power to a power \textbf{multiplies} the exponents
($\left(2^{3}\right)^{x}=2^{3x}$, never $2^{\,3+x}$), and when the exponent is a binomial the
multiplication has to reach \emph{both} terms ($\left(3^{2}\right)^{x+4}=3^{\,2x+8}$). Step 2 is legal
because of a fact you proved back in Lesson 7.0: an exponential graph has \textbf{no turning points},
so it never uses the same output twice --- which means $b^{m}=b^{n}$ can only happen when $m=n$. That
is also why the base has to satisfy $b>0$ and $b\neq1$, the restriction Lesson 7.1 handed you: $1^{5}$
and $1^{9}$ are both $1$, so with $b=1$ matching bases would prove nothing. Finally, keep two kinds of
failure apart. $2^{x}=-3$ has \textbf{no solution}, because $-3$ is not in the range $(0,\infty)$.
$2^{x}=5$ \textbf{does} have a solution --- the line $y=5$ crosses the curve between $x=2$ and $x=3$
--- but $5$ is not a power of $2$, so today's method cannot reach it. That second kind is what Lesson
7.4 walks up to and Unit 8 finally solves.
\end{remindbox}

\end{document}
